% Clinical Background section. Paste into your own thesis, or \input it.

\section{Geometry as a risk factor}
\label{sec:geometry}

The preceding sections describe the coronary tree and the disease that affects it.
This section states why the \emph{shape} of the tree is a clinical quantity rather than an
incidental descriptor.

The accepted causal chain runs from geometry through hemodynamics to biology: geometry and
topology determine wall shear stress, wall shear stress drives the endothelial response, and
the endothelial response produces plaque \cite{rampidis2022geometry}.

Wall shear stress (WSS) is the tangential frictional force exerted by flowing blood on the
endothelium.
Regions of low and oscillatory shear are where early atherosclerotic lesions preferentially
develop, because the endothelium responds to disturbed shear with a pro-inflammatory
phenotype \cite{rampidis2022geometry}.
The geometric configurations that produce low shear are predictable.
Flow separates at a bifurcation, leaving a low-shear region on the lateral walls away from
the flow divider, and curvature drives secondary flows that leave a low-shear region on the
inner wall of a bend \cite{rampidis2022geometry}.
Bifurcation angle and tortuosity are therefore not arbitrary descriptors of shape; they are
proxies for the size and location of the regions in which lesions form.

This reasoning is supported by outcome data.
In a study of patients who subsequently experienced an acute coronary syndrome, three
geometric characteristics carried prognostic value beyond stenosis severity and plaque
characteristics: a more proximal lesion location, location at a bifurcation, and increased
tortuosity \cite{han2022plaque}.
That result is the strongest available justification for treating the geometry of the
coronary tree as an object of study in its own right.

One quantitative anchor is worth recording: the normal left main bifurcation angle has been
reported as $76.4^\circ \pm 16.7^\circ$ \cite{rampidis2022geometry}.
Curvature has been reported as higher in atherosclerotic than in non-atherosclerotic segments,
and higher in end-systole than in diastole \cite{rampidis2022geometry}.
The direction of the association between bifurcation angle and risk is not consistent across
the literature, so its sign should not be assumed in advance.

A scaling relation completes the picture.
Murray's law, derived by minimizing the sum of the power dissipated in flow and the metabolic
cost of maintaining blood, predicts that flow scales with the cube of vessel radius, and
hence that wall shear stress is constant throughout a healthy arterial tree.
In the coronary circulation specifically the exponent is not three.
A meta-analysis over 1070 coronary trees places it at 2.39, with a 95\% confidence interval
of 2.24 to 2.54 \cite{taylor2024murray}.
Any feature defined as a deviation from Murray's law must therefore use the coronary
exponent rather than the cubic one.

These ideas have been applied to coronary morphology directly.
Wang et al.\ extracted centerline geometry from 76 left main arteries, applied unsupervised
clustering to obtain four morphological phenotypes, and computed time-averaged wall shear
stress for each by computational fluid dynamics, finding low shear around the branch points of
the LAD and most markedly in the phenotype characterized by a large bifurcation angle
\cite{wang2024lmphenotypes}.
